% ==============================================================================
% Capítulo 1: Introdução
% ==============================================================================
\chapter{Introdução}
\label{cap:introducao}

\begin{instrucao}
\orientacao{O que vai na Introdução}
A Introdução situa quem lê o trabalho. Ela apresenta o contexto do tema, o problema
tratado, os objetivos, a justificativa e a organização do texto. Como
referência, de três a cinco páginas costumam ser suficientes. Escreva uma
primeira versão no início do trabalho e reescreva no final, quando os
resultados estiverem claros, porque a Introdução promete o que o resto do texto
precisa cumprir.

Comece pelo geral e avance até o problema específico. Um bom primeiro parágrafo
responde ``por que este assunto importa?'' sem frases genéricas como ``com o
avanço da tecnologia\ldots''.

\orientacao{Capítulo ou seção?}
A NBR 14724:2024 determina que o trabalho seja dividido em seções, e não em
capítulos. No LaTeX nada muda (continue usando \verb|\chapter|, que gera o
título ``1 INTRODUÇÃO''). A diferença está só no nome usado nas referências
cruzadas, que este modelo controla pelo comando \verb|\termocap| em
\verb|config/estilo.tex|. O padrão do modelo é ``Capítulo'', o termo mais
usado nos cursos.

\orientacao{Rótulos}
Logo após cada \verb|\chapter| ou \verb|\section|, coloque um
\verb|\label| com um prefixo que indique o tipo: \verb|cap:| para capítulos,
\verb|sec:| para seções, \verb|fig:| para figuras, \verb|tab:| para tabelas,
\verb|qua:| para quadros, \verb|eq:| para equações e assim por diante. Nunca
use acentos ou espaços nos rótulos.
\end{instrucao}

\begin{instrucao}
\orientacao{Pessoa verbal e estilo acadêmico}
Textos acadêmicos são escritos, por convenção, de forma impessoal. A NBR
14724:2024 não determina a pessoa verbal do trabalho, mas a NBR 6028:2021
recomenda o verbo na terceira pessoa no resumo (item 4.1.4). Este modelo
estende essa recomendação a todo o texto, para que ele fique uniforme. Use:
\begin{itemize}
  \item a terceira pessoa do singular, com o trabalho, o método ou o sistema
        como sujeito (``este trabalho propõe'', ``a heurística obteve'',
        ``o sistema permite'');
  \item ou a voz passiva (``os experimentos foram executados'', ``observa-se
        que'', ``adotou-se'').
\end{itemize}
Evite a primeira pessoa do singular (``eu proponho'') e do plural (``nós
propomos'', ``vamos mostrar''), mesmo como plural de modéstia. Algumas áreas
aceitam o ``nós'', mas a forma impessoal é a convenção predominante nos
trabalhos acadêmicos. Qualquer que seja a escolha, não misture as duas formas
no mesmo texto, um erro frequente nos TCCs. As exceções são
os agradecimentos, a dedicatória e a declaração de uso de inteligência
artificial, escritos na primeira pessoa do singular por serem manifestações
pessoais do autor.

\begin{itemize}
  \item \textbf{Evite:} ``Neste trabalho, vamos propor uma heurística e
        mostraremos que ela é melhor.''
  \item \textbf{Prefira:} ``Este trabalho propõe uma heurística e avalia seu
        desempenho em relação a \ldots''
\end{itemize}

Algumas regras de redação ajudam o texto a soar claro e profissional:
\begin{itemize}
  \item escreva frases diretas e evite expressões de preenchimento como ``vale
        ressaltar que'' e ``é importante notar que'';
  \item evite adjetivos de efeito (``crucial'', ``poderoso'', ``robusto'') e
        afirmações sem dados ou citação (``muito usado'', ``o melhor método'');
  \item não use linguagem coloquial nem perguntas retóricas;
  \item use vírgula decimal e ponto como separador de milhar (0,41\% e
        1.200), inclusive em tabelas e gráficos.
\end{itemize}
\end{instrucao}

Muitas decisões de planejamento envolvem escolher, dentre várias alternativas,
aquelas que trazem o maior retorno sem ultrapassar recursos limitados. Uma
universidade que seleciona projetos de extensão para financiar precisa
respeitar ao mesmo tempo o orçamento, a carga horária dos docentes e o número
de bolsas disponíveis. Um provedor de computação em nuvem que aloca tarefas a
um servidor precisa respeitar a capacidade de processamento, de memória e de
armazenamento da máquina. Em ambos os casos, cada alternativa consome uma
quantidade diferente de cada recurso.

Esse tipo de decisão é modelado pelo \ac{mkp}, um dos problemas mais estudados
da otimização combinatória \cite{freville2004multidimensional,
puchinger2010multidimensional}. No \acs{mkp}, cada item tem um valor e um peso
em cada uma de várias dimensões e o objetivo é escolher o subconjunto de itens
de maior valor total que respeite a capacidade de todas as dimensões.

\begin{instrucao}
\orientacao{Siglas}
Declare cada sigla uma única vez no arquivo \verb|cientifico/siglas.tex| e use
\verb|\ac{chave}| no texto. Na primeira ocorrência o LaTeX escreve a forma
extensa seguida da sigla. Nas seguintes, apenas a sigla. A lista de siglas é
montada sozinha, só com as siglas usadas.
\begin{exemplo}
A \ac{abnt} publica as normas.
Este trabalho segue as
normas da \ac{abnt}.
\end{exemplo}
\acreset{abnt}% desfaz o uso da sigla no exemplo acima
Em títulos e legendas use \verb|\acs{chave}|, que imprime só a sigla e não
interfere na contagem da primeira ocorrência. Não use \verb|\ac| no resumo, que
deve ser compreensível sem consultar a lista.
\end{instrucao}

% ------------------------------------------------------------------------------
\section{Problema e questão de pesquisa}
\label{sec:introducao:problema}

O \acs{mkp} pertence à classe dos problemas NP-difíceis \cite{kellerer2004knapsack}.
Na prática, isso significa que não se conhece algoritmo capaz de resolver todas
as instâncias de forma exata em tempo polinomial. Solucionadores de \ac{pli}
encontram a solução ótima de instâncias pequenas em poucos segundos, mas o
tempo de execução cresce rapidamente com o número de itens e de restrições.
Para instâncias com 500 itens e 30 restrições, mesmo solucionadores comerciais
podem levar horas sem provar a otimalidade \cite{puchinger2010multidimensional}.

Nesses casos, heurísticas e metaheurísticas são a alternativa mais usada,
porque encontram soluções de boa qualidade em tempo limitado, ainda que sem
garantia de otimalidade. Diante disso, este trabalho busca responder à
seguinte questão de pesquisa: \emph{uma heurística de \ac{ils} com perturbação
adaptativa consegue obter, em até 60 segundos, soluções para o \acs{mkp} mais
próximas das melhores conhecidas do que um solucionador de \acs{pli} no mesmo
tempo?}

\begin{instrucao}
\orientacao{Problema, questão de pesquisa e hipótese}
Em um trabalho científico, deixe explícito qual pergunta o trabalho responde.
Uma boa questão de pesquisa é específica e verificável pelos experimentos. A
pergunta acima define o método (\acs{ils}), a comparação (solucionador de
\acs{pli}), a métrica (proximidade das melhores soluções) e a condição (60
segundos). Se preferir, formule uma hipótese, que é uma resposta provisória à
pergunta e será confirmada ou refutada pelos resultados.
\end{instrucao}

% ------------------------------------------------------------------------------
\section{Objetivos}
\label{sec:introducao:objetivos}

\begin{instrucao}
\orientacao{Como escrever os objetivos}
O \textbf{objetivo geral} expressa o resultado principal do trabalho em uma
única frase iniciada por verbo no infinitivo. Os \textbf{objetivos
específicos} são os resultados intermediários que, juntos, levam ao objetivo
geral. Cada um deve ser verificável, ou seja, ao final do trabalho deve ser
possível dizer se foi alcançado e em que capítulo isso aparece.

Prefira verbos que indiquem um resultado concreto: \emph{propor},
\emph{implementar}, \emph{comparar}, \emph{avaliar}, \emph{identificar},
\emph{desenvolver}, \emph{analisar}. Evite verbos vagos como \emph{estudar},
\emph{conhecer}, \emph{entender} ou \emph{aprender}, que descrevem o processo
de estudo, e não o resultado do trabalho.

\begin{itemize}
  \item \textbf{Fraco:} ``Estudar metaheurísticas para o problema da mochila.''
  \item \textbf{Melhor:} ``Propor e avaliar uma heurística de busca local
        iterada para o problema da mochila multidimensional.''
\end{itemize}

Na Conclusão, retome cada objetivo específico e diga como ele foi atingido.
Para listas numeradas por letras, como pede a NBR 6024 para alíneas, use o
ambiente \verb|alineas| do abnTeX2.
\begin{exemplo}
\begin{alineas}
  \item primeiro item;
  \item segundo item.
\end{alineas}
\end{exemplo}
\end{instrucao}

O objetivo geral deste trabalho é propor e avaliar uma heurística de \ac{ils}
para o \acs{mkp} capaz de obter soluções de boa qualidade em tempo limitado.
Para alcançá-lo, foram definidos os seguintes objetivos específicos:
\begin{alineas}
  \item identificar, na literatura, as principais abordagens heurísticas para
        o \acs{mkp} e as instâncias usadas para avaliá-las;
  \item propor uma heurística de \acs{ils} com perturbação adaptativa para o
        problema;
  \item implementar a heurística e um modelo de \acs{pli} resolvido pelo
        solucionador HiGHS;
  \item comparar o desempenho da heurística com o do solucionador e com o de
        um \ac{ag} da literatura em instâncias de referência;
  \item avaliar a contribuição de cada componente da heurística para o
        resultado final.
\end{alineas}

% ------------------------------------------------------------------------------
\section{Justificativa}
\label{sec:introducao:justificativa}

\begin{instrucao}
\orientacao{Justificativa e citações}
A justificativa mostra por que o trabalho é relevante. Apoie-se em trabalhos
anteriores, aponte as limitações deles e mostre a lacuna que o seu trabalho
procura preencher. Toda afirmação que não seja de conhecimento geral da área
precisa de citação.

Este modelo usa o biblatex-abnt, que segue a NBR 10520:2023. Na nova norma, o
sobrenome do autor aparece com apenas a inicial maiúscula também dentro dos
parênteses, como em (Chu; Beasley, 1998). Use \verb|\textcite| quando o autor
faz parte da frase e \verb|\cite| quando a citação vai entre parênteses.
\begin{exemplo}
Segundo \textcite{chu1998genetic},
o problema é difícil.

O problema é difícil
\cite{chu1998genetic}.

Com página:
\cite[p.~70]{chu1998genetic}.

Várias obras: \cite{chu1998genetic,
kellerer2004knapsack}.
\end{exemplo}
As chaves (como \verb|chu1998genetic|) são definidas no arquivo
\verb|referencias.bib|. Só entram nas Referências as obras citadas no texto.
Para citação de citação (\emph{apud}), use
\verb|\apud{obra-original}{obra-consultada}|, lembrando que a obra original
não entra nas referências se você não teve acesso a ela.
\end{instrucao}

O \acs{mkp} tem sido estudado há décadas e conta com um conjunto de instâncias
de referência amplamente usado, disponibilizado por \textcite{chu1998genetic}
na OR-Library\footnote{Disponível em
\url{https://people.brunel.ac.uk/~mastjjb/jeb/orlib/mknapinfo.html}.}. Esse
conjunto permite comparar novas abordagens com resultados já publicados, o que
torna o problema adequado para avaliar ideias de projeto de heurísticas.

Apesar do grande número de métodos propostos, muitos dependem de
procedimentos sofisticados, como a combinação de programação linear com busca
tabu \cite{vasquez2001hybrid}, que são difíceis de reproduzir e de adaptar a
variantes do problema \cite{puchinger2010multidimensional}. Uma heurística
simples, com poucos parâmetros e bom desempenho em tempo limitado, é útil
tanto como método de solução quanto como ponto de partida para variantes com
restrições adicionais.

% ------------------------------------------------------------------------------
\section{Organização do trabalho}
\label{sec:introducao:organizacao}

\begin{instrucao}
\orientacao{Organização do texto}
Feche a Introdução descrevendo, em um parágrafo ou em uma lista curta, o que
cada capítulo apresenta. Use \verb|\Cref{rotulo}| no início da frase e
\verb|\cref{rotulo}| no meio dela. O número é atualizado sozinho se a ordem dos
capítulos mudar.
\end{instrucao}

Além desta introdução, o trabalho está organizado da seguinte forma. O
\cref{cap:fundamentacao} apresenta a formulação matemática do \acs{mkp} e os
conceitos de metaheurísticas usados no trabalho. O \cref{cap:relacionados}
discute os principais trabalhos da literatura sobre o problema. O
\cref{cap:metodologia} descreve a heurística proposta e o planejamento dos
experimentos, cujos resultados são apresentados e discutidos no
\cref{cap:experimentos}. Por fim, o \cref{cap:conclusao} retoma os objetivos,
resume as contribuições e aponta trabalhos futuros.
